\documentclass{article}
\usepackage{annot}
\usetikzlibrary{decorations.pathreplacing}
% Accolade à main levée : \accolade[couleur]{(x1,y)}{(x2,y)} (dessous si x1 > x2)
\NewDocumentCommand{\accolade}{O{violet} m m}{%
  \draw[crayon lisse, #1, decorate, decoration={brace, amplitude=1.4mm}] #2 -- #3;}
% Cadre à main levée (évite \encadrer, qui déborde sur les grands rectangles)
\NewDocumentCommand{\cadre}{O{rouge} m m m m}{%
  \draw[crayon, sharp corners, #1] (#2,#3) -- (#4,#3) -- (#4,#5) -- (#2,#5) -- cycle;}
\begin{document}

% ---------------------------------------------------------------------------
% Notion d'intervalle de confiance
% ---------------------------------------------------------------------------
\annot{3}{Limite de l'estimation ponctuelle}{
  \ecrire[violet][8]{(38,52)}{$\overline{X}$ : VARIABLE ALÉATOIRE \quad $\neq$ \quad $\bar{x} = 134$ : UN NOMBRE (RÉALISATION)}
  \entourer[rouge]{(101,25)}{4.5}{2.6}
}

\annot{4}{Définition d'un intervalle de confiance}{
  \ecrire[violet][9]{(62,81.5)}{EX. : IC À 95\% $\Rightarrow 1-\alpha = 0{,}95$, \ $\alpha = 0{,}05$}
  \surligner[jaune]{(71,50.5)}{(106,56.5)}
  % petit schéma de l'intervalle
  \trait[vert]{(124,13) -- (150,13)}
  \trait[vert]{(124,11.3) -- (124,14.7)}
  \trait[vert]{(150,11.3) -- (150,14.7)}
  \fill[vert] (137,13) circle (0.8);
  \ecrirec[vert][7]{(124,9.3)}{$L_1$}
  \ecrirec[vert][7]{(150,9.3)}{$L_2$}
  \ecrirec[vert][6]{(137,9.3)}{ESTIM.}
  \accolade[rouge]{(124.5,15)}{(136.5,15)}
  \accolade[rouge]{(137.5,15)}{(149.5,15)}
  \ecrirec[rouge][7]{(130.5,19.2)}{ME}
  \ecrirec[rouge][7]{(143.5,19.2)}{ME}
}

\annot{5}{Que veut dire}{
  \ecrire[rouge][8]{(47,82)}{$\theta$ : FIXE (MAIS INCONNU)}
  \fleche[rouge][-25]{(52,77.5)}{(41.8,71.5)}
  \croix{(6.5,53)}
  \croix{(64.5,14.5)}
  \ecrire[rouge][7]{(2,49.5)}{RATE $\theta$}
  \ecrire[violet][8.5]{(10,8.5)}{CE SONT LES INTERVALLES QUI SONT ALÉATOIRES !}
}

\annot{6}{Simulation : 40 intervalles}{
  \entourer[rouge]{(36.2,50.5)}{2.2}{4.3}
  \entourer[rouge]{(70.2,40)}{2.2}{4.3}
  \entourer[rouge]{(84.3,39.5)}{2.2}{5.5}
  \entourer[rouge]{(138.2,40.3)}{2.2}{4.3}
  \ecrire[rouge][7]{(31,60.5)}{RATE $\mu$ !}
  \ecrire[violet][9]{(10,13)}{CHAQUE BARRE = 1 ÉCHANTILLON DE 7 PERSONNES\\
    $\rightarrow$ SON PROPRE CENTRE $\bar{x}$ ET SA PROPRE LONGUEUR ($s$)}
}

\annot{7}{Interprétation rigoureuse}{
  \ecrire[rouge][9]{(12,82.5)}{UNE FOIS CALCULÉ, L'IC CONTIENT $\mu$ OU NON : \  PROBA 0 OU 1 !}
  \surligner[jaune]{(95,53.5)}{(150,58.5)}
  \surligner[jaune]{(10,48.5)}{(47,53.5)}
  \ecrire[violet][7.5]{(51,52.8)}{$\leftarrow$ C'EST LA MÉTHODE QUI EST FIABLE À 95\%}
}

\annot{8}{QUIZ}{
  \ecrire[vert][10]{(18,66)}{$40 \times 0{,}10 = 4$ INTERVALLES EN MOYENNE\ \ (ICI : EXACTEMENT 4)}
  \ecrire[vert][9]{(18,45.5)}{$40 \times 0{,}05 = 2$ EN MOYENNE ; \ PLUS LONGS :\\
    QUANTILE PLUS GRAND \ ($t_{6;0{,}025} = 2{,}447 > t_{6;0{,}05} = 1{,}943$)}
  \ecrire[vert][10]{(18,24)}{CAR $s$ CHANGE D'UN ÉCHANTILLON À L'AUTRE :\\
    LONGUEUR $= 2\, t_{6;0{,}05}\, \dfrac{s}{\sqrt{7}}$ \ ($t$ ET $n$ FIXES, $s$ ALÉATOIRE)}
}

% ---------------------------------------------------------------------------
% Construction de l'intervalle pour la moyenne
% ---------------------------------------------------------------------------
\apres{9}{
  \ecrire[violet][15]{(7,85)}{SOIT \ $Z \sim N(0,1)$}
  \ecrire[violet][15]{(7,74)}{$P(Z < 1) = 0{,}8413$}
  \trait[orange]{(7,62.5) -- (72,62)}
  \ecrire[orange][12]{(7,59)}{$P(Z \le k) = 0{,}8413 \ \Rightarrow\ k = 1$}
  \ecrire[orange][12]{(7,50)}{$P(Z \le w) = 0{,}9332 \ \Rightarrow\ w = 1{,}5$}
  \ecrire[vert][12]{(7,39)}{JE VEUX $a$ ET $b$ T.Q.\\[1mm] $P(a \le Z \le b) = 0{,}90$}
  \ecrire[vert][12]{(7,21)}{POUR UN IC : $a = -b$\\[1mm] $\Rightarrow P(Z > b) = 0{,}05$\\[1mm] $\Rightarrow b = z_{0{,}05} = 1{,}645$}
  \ecrire[violet][11]{(88,86)}{POUR UN IC DE NIVEAU $1-\alpha$ :\\[1mm]
    $P(-z_{\alpha/2} \le Z \le z_{\alpha/2}) = 1-\alpha$}
  % courbe normale
  \begin{scope}[shift={(123,13)}]
    \fill[pattern=north east lines, pattern color=vert] (-15.63,0) --
      plot[domain=-1.645:1.645, samples=50] ({9.5*\x},{38*exp(-(\x)^2/2)}) -- (15.63,0) -- cycle;
    \fill[rose, opacity=.55] (-30.4,0) --
      plot[domain=-3.2:-1.645, samples=30] ({9.5*\x},{38*exp(-(\x)^2/2)}) -- (-15.63,0) -- cycle;
    \fill[rose, opacity=.55] (15.63,0) --
      plot[domain=1.645:3.2, samples=30] ({9.5*\x},{38*exp(-(\x)^2/2)}) -- (30.4,0) -- cycle;
    \draw[crayon lisse, orange, line width=1.1pt] plot[domain=-3.2:3.2, samples=90, smooth] ({9.5*\x},{38*exp(-(\x)^2/2)});
    \trait[violet]{(-31.5,0) -- (31.5,0)}
    \trait[orange]{(0,0) -- (0,38)}
    \ecrirec[violet][11]{(-15.6,-4)}{$-b$}
    \ecrirec[violet][11]{(0,-4)}{$0$}
    \ecrirec[violet][11]{(15.6,-4)}{$b$}
    \ecrirec[vert][12]{(-20,27)}{$1-\alpha$}
    \fleche[vert][-30]{(-15,24)}{(-4,14)}
    \ecrirec[rose][11]{(-27,13)}{$\alpha/2$}
    \fleche[rose][25]{(-27,9.5)}{(-21,2)}
    \ecrirec[rose][11]{(27,13)}{$\alpha/2$}
    \fleche[rose][-25]{(27,9.5)}{(21,2)}
    \ecrire[orange][10]{(8,40)}{DENSITÉ DE $Z$}
  \end{scope}
}

\annot{10}{Cas théorique}{
  \ecrire[violet][8]{(122,66)}{$\leftarrow$ CENTRER-\\ \ \ \ RÉDUIRE}
  \surligner[jaune]{(103.5,14.8)}{(122.5,19.2)}
  \ecrirec[vert][7]{(84,12.3)}{90\%}
  \ecrirec[vert][7]{(113,12.3)}{95\%}
  \ecrirec[vert][7]{(136,12.3)}{99\%}
  \ecrire[rouge][8]{(126,10)}{$\Rightarrow$ ON UTILISE $S$}
}
\apres{10}{
  \ecrire[violet][13]{(7,86)}{ON ISOLE $\mu$ AU CENTRE :}
  \ecrire[vert][11]{(7,76)}{$P\left(-z_{\alpha/2} \le \dfrac{\overline{X} - \mu}{\sigma/\sqrt{n}} \le z_{\alpha/2}\right) = 1-\alpha$}
  \ecrire[vert][11]{(7,60)}{$P\left(-z_{\alpha/2}\dfrac{\sigma}{\sqrt{n}} \le \overline{X} - \mu \le z_{\alpha/2}\dfrac{\sigma}{\sqrt{n}}\right) = 1-\alpha$}
  \ecrire[vert][11]{(7,44)}{$P\left(-\overline{X} - z_{\alpha/2}\dfrac{\sigma}{\sqrt{n}} \le -\mu \le -\overline{X} + z_{\alpha/2}\dfrac{\sigma}{\sqrt{n}}\right) = 1-\alpha$}
  \ecrire[vert][11]{(7,28)}{$P\left(\overline{X} - z_{\alpha/2}\dfrac{\sigma}{\sqrt{n}} \le \mu \le \overline{X} + z_{\alpha/2}\dfrac{\sigma}{\sqrt{n}}\right) = 1-\alpha$}
  \cadre[rouge]{5}{15.5}{112}{29.5}
  \ecrire[violet][10]{(122,56)}{$\times\ \dfrac{\sigma}{\sqrt{n}}$}
  \ecrire[violet][10]{(128,40)}{$-\ \overline{X}$}
  \ecrire[rouge][10]{(116,27)}{$\times (-1)$ :\\ LES $\le$ CHANGENT\\ DE SENS !}
  \accolade[orange]{(38,15)}{(16,15)}
  \accolade[orange]{(89,15)}{(66,15)}
  \ecrirec[orange][11]{(27,8.5)}{$L_1$}
  \ecrirec[orange][11]{(77.5,8.5)}{$L_2$}
  \ecrire[violet][9]{(96,12)}{$L_1$, $L_2$ ALÉATOIRES\\ (DÉPENDENT DE $\overline{X}$)}
}

\annot{11}{Que faire quand}{
  \entourer[rouge]{(80.3,58.2)}{2.5}{2.5}
  \ecrire[rouge][9]{(98,70)}{$\sigma$ INCONNUE $\rightarrow$ ESTIMÉE PAR $S$}
  \fleche[rouge][25]{(100,65.5)}{(83.2,59.5)}
  \ecrire[violet][8.5]{(72,40)}{$\rightarrow$ QUANTILES PLUS GRANDS QUE 1,96}
}

\annot{12}{William Sealy Gosset}{
  \entourer[orange]{(24,66)}{7.5}{2.6}
  \entourer[orange]{(48,21)}{9.5}{2.8}
  \ecrire[violet][9]{(12,11)}{$n$ PETIT $\Rightarrow$ $S$ PEU PRÉCIS $\Rightarrow$ IL FAUT LA LOI EXACTE DE $T$}
}

\annot{13}{La loi de Student}{
  \ecrire[violet][8]{(56,83.5)}{$S^2 = \frac{1}{n-1}\sum (X_i - \overline{X})^2$ : ON PERD 1 D.L. (ON ESTIME $\mu$ PAR $\overline{X}$)}
  \ecrire[rouge][8]{(121,60)}{DEGRÉS DE\\ LIBERTÉ}
  \fleche[rouge][25]{(121,54)}{(107,48.8)}
  \ecrire[vert][8]{(100,43.5)}{EX. : $n = 7 \Rightarrow k = 6$}
}

\annot{14}{Comparaison de la loi normale}{
  \ecrire[rouge][7.5]{(41.5,72)}{$k \rightarrow \infty$ :\\ STUDENT\\ $\rightarrow N(0,1)$}
  \ecrire[rouge][8]{(131,55)}{QUEUES PLUS\\ LOURDES}
  \fleche[rouge][-20]{(135,47)}{(125,37.2)}
  \entourer[vert]{(47.5,13.8)}{5.2}{2.4}
  \entourer[vert]{(76,13.8)}{5.2}{2.4}
  \ecrire[vert][8.5]{(84,14)}{$>$ \ $z_{0{,}025} = 1{,}96$ \ (NORMALE)}
}

\annot{15}{Formule de l'IC pour}{
  \ecrire[rouge][8]{(12,56)}{DEGRÉS DE\\ LIBERTÉ}
  \fleche[rouge][-20]{(22,49.5)}{(68.5,46.6)}
  \accolade[violet]{(122.5,44.2)}{(93.5,44.2)}
  \ecrire[violet][8]{(124,51)}{MARGE\\ D'ERREUR}
  \ecrire[violet][8]{(73,30.5)}{(ÉCART-TYPE DE $\overline{X}$, ESTIMÉ)}
}
\apres{15}{
  \ecrire[violet][14]{(7,86)}{TABLE DE STUDENT}
  \ecrire[violet][10]{(7,76.5)}{$P(T_k > t_{k;\gamma}) = \gamma$}
  % petite courbe avec queue droite hachurée
  \begin{scope}[shift={(90,67)}]
    \fill[pattern=north east lines, pattern color=rouge] (8,0) --
      plot[domain=1.6:3, samples=20] ({5*\x},{14*exp(-(\x)^2/2)}) -- (15,0) -- cycle;
    \draw[crayon lisse, rouge] plot[domain=-3:3, samples=60, smooth] ({5*\x},{14*exp(-(\x)^2/2)});
    \trait[rouge]{(-16,0) -- (16,0)}
    \ecrirec[rouge][8]{(21,8)}{$\gamma$}
    \fleche[rouge][-20]{(20,5.5)}{(12,1.5)}
    \ecrirec[rouge][8]{(8,-3)}{$t_{k;\gamma}$}
  \end{scope}
  % la table
  \begin{scope}[bleu]
    \ecrirec[bleu][10]{(12,64)}{$k$}
    \foreach \g [count=\j] in {{0,10},{0,05},{0,025},{0,01},{0,005}}{
      \ecrirec[bleu][10]{({14+17*\j},64)}{\g}}
    \trait[bleu]{(5,60) -- (112,60)}
    \trait[bleu]{(19,68) -- (19,3)}
  \end{scope}
  \foreach \k/\a/\b/\c/\d/\e [count=\i] in {
      4/1{,}533/2{,}132/2{,}776/3{,}747/4{,}604,
      5/1{,}476/2{,}015/2{,}571/3{,}365/4{,}032,
      6/1{,}440/1{,}943/2{,}447/3{,}143/3{,}707,
      7/1{,}415/1{,}895/2{,}365/2{,}998/3{,}499,
      8/1{,}397/1{,}860/2{,}306/2{,}896/3{,}355}{
    \ecrirec[bleu][10]{(12,{63-7*\i})}{$\k$}
    \ecrirec[bleu][10]{(31,{63-7*\i})}{$\a$}
    \ecrirec[bleu][10]{(48,{63-7*\i})}{$\b$}
    \ecrirec[bleu][10]{(65,{63-7*\i})}{$\c$}
    \ecrirec[bleu][10]{(82,{63-7*\i})}{$\d$}
    \ecrirec[bleu][10]{(99,{63-7*\i})}{$\e$}}
  \ecrirec[bleu][10]{(12,21.5)}{$\vdots$}
  \ecrirec[bleu][10]{(56,21.5)}{$\vdots$}
  \foreach \k/\a/\b/\c/\d/\e [count=\i] in {
      29/1{,}311/1{,}699/2{,}045/2{,}462/2{,}756,
      \infty/1{,}282/1{,}645/1{,}960/2{,}326/2{,}576}{
    \ecrirec[bleu][10]{(12,{20-7*\i})}{$\k$}
    \ecrirec[bleu][10]{(31,{20-7*\i})}{$\a$}
    \ecrirec[bleu][10]{(48,{20-7*\i})}{$\b$}
    \ecrirec[bleu][10]{(65,{20-7*\i})}{$\c$}
    \ecrirec[bleu][10]{(82,{20-7*\i})}{$\d$}
    \ecrirec[bleu][10]{(99,{20-7*\i})}{$\e$}}
  \surligner[rose]{(6,38.3)}{(109,43.7)}
  \surligner[jaune]{(41.5,3)}{(54.5,67)}
  \entourer[vert]{(48,41)}{6}{3.2}
  \entourer[orange]{(65,41)}{6.3}{3.2}
  \ecrire[vert][11]{(116,57)}{IC À 90\% :\\ $\alpha = 0{,}10$\\ $\alpha/2 = 0{,}05$\\ $n = 7 \Rightarrow k = 6$\\ $\Rightarrow t_{6;0{,}05} = 1{,}943$}
  \ecrire[orange][9]{(116,26)}{IC À 95\% :\\ $t_{6;0{,}025} = 2{,}447$}
  \ecrire[rouge][9]{(116,13.5)}{$k = \infty$ : ON\\ RETROUVE $z$ !}
  \fleche[rouge][-20]{(115.5,9.5)}{(108,6)}
}

\annot{16}{Exemple 1 (suite) : pression}{
  \ecrire[violet][8]{(52,70.5)}{$s^2 = \frac{1}{6}\sum (x_i - 134)^2 = \frac{236}{6} = 39{,}33$}
  \ecrire[rouge][8]{(101,56.3)}{$\leftarrow$ DEGRÉS DE LIBERTÉ}
  \surligner[jaune]{(107.5,45.4)}{(118,49.6)}
  \ecrire[violet][7]{(123,50)}{$\leftarrow$ TABLE : LIGNE 6,\\ \ \ \ COLONNE 0,05}
  % schéma de l'IC
  \trait[vert]{(112,21) -- (148,21)}
  \trait[vert]{(112,19.2) -- (112,22.8)}
  \trait[vert]{(148,19.2) -- (148,22.8)}
  \fill[vert] (130,21) circle (0.8);
  \ecrirec[vert][6.5]{(112,17.3)}{129,39}
  \ecrirec[vert][6.5]{(130,17.3)}{134}
  \ecrirec[vert][6.5]{(148,17.3)}{138,61}
  \accolade[rouge]{(112.5,23.2)}{(129.5,23.2)}
  \accolade[rouge]{(130.5,23.2)}{(147.5,23.2)}
  \ecrirec[rouge][7]{(121,27.6)}{4,61}
  \ecrirec[rouge][7]{(139,27.6)}{4,61}
}

\annot{17}{Avec R}{
  \entourer[vert]{(49,60.4)}{29}{2.5}
  \ecrire[vert][8.5]{(84,62.8)}{$= [\,129{,}39 \,;\, 138{,}61]$ \ COMME À LA MAIN}
  \ecrire[vert][8]{(84,44)}{$\leftarrow$ 95\% : PLUS LONG ($2{,}447 > 1{,}943$)}
  \ecrire[violet][8]{(54,14)}{$0{,}95 = 1-\alpha/2$ : AIRE À GAUCHE}
  \fleche[violet][-20]{(58,14.2)}{(65,16.5)}
}

% ---------------------------------------------------------------------------
% Quiz conceptuel : vrai ou faux
% ---------------------------------------------------------------------------
\annot{19}{Question 1}{
  \entourer[violet]{(18.5,42.8)}{10}{3.8}
  \ecrire[violet][9.5]{(10,15.5)}{L'IC PORTE SUR $\mu$, PAS SUR LES $X$ INDIVIDUELS}
  \ecrire[vert][8.5]{(10,8.5)}{90\% DES INDIVIDUS : ENVIRON $134 \pm 1{,}645 \times 6{,}27 = [\,123{,}7 \,;\, 144{,}3]$ \ (BEAUCOUP PLUS LARGE !)}
}

\annot{20}{Question 2}{
  % axe des pressions : x = 35 + (v - 120) * 4.4
  \trait[violet]{(35,41) -- (148,41)}
  \foreach \v in {130,135,137,138}{\fill[vert] ({35+(\v-120)*4.4},42.2) circle (0.8);}
  \fill[vert] ({35+15*4.4},44.2) circle (0.8);
  \croix{({35+2*4.4},42.4)}
  \croix{({35+21*4.4},42.4)}
  \ecrirec[rouge][7]{({35+2*4.4},46.8)}{122}
  \ecrirec[rouge][7]{({35+21*4.4},46.8)}{141}
  \draw[crayon lisse, rouge, line width=1pt] (76.3,45) -- (76.3,47.3) -- (116.9,47.3) -- (116.9,45);
  \ecrirec[rouge][7]{(96.5,50.2)}{IC À 90\%}
  \ecrire[vert][8]{(118,33.5)}{5 SUR 7 SEULEMENT}
}

\annot{21}{Question 3}{
  \ecrire[violet][8]{(30,46)}{$\overline{X}_{\text{NOUV}} \sim N\!\left(\mu, \frac{\sigma^2}{n}\right)$ : VARIE AUTOUR DE $\mu$,\\ ET L'IC (CENTRÉ SUR $\bar{x}$) VARIE AUSSI}
  \cloche[orange]{(134,36.5)}{28}{9}
  \trait[orange]{(134,36.5) -- (134,45.5)}
  \ecrire[orange][7]{(135,45)}{$\mu$}
  \ecrire[rouge][8.5]{(60,8.5)}{$\Rightarrow$ PROBA $\neq 90\%$ (EN MOYENNE, NETTEMENT MOINS)}
}

\annot{22}{Question 4}{
  \ecrire[vert][10]{(36,42.5)}{$1000 \times 0{,}90 = 900$ \ (ENVIRON)}
  \ecrire[vert][10]{(10,16.5)}{CAR \ $P\left(\overline{X} - t_{6;0{,}05}\dfrac{S}{\sqrt{7}} \le \mu \le \overline{X} + t_{6;0{,}05}\dfrac{S}{\sqrt{7}}\right) = 0{,}90$}
  \ecrire[violet][8]{(104,12.5)}{($\mu$ FIXE,\\ BORNES ALÉATOIRES)}
}

\annot{23}{Question 5}{
  \entourer[vert]{(42.5,23.9)}{13.5}{2.8}
  \ecrire[violet][9]{(10,11.5)}{PAR CONTRE : POPULATION NORMALE ET $\sigma$ INCONNUE\\ $\Rightarrow$ ON UTILISE LA STUDENT (UN PEU PLUS LONG = PRUDENT)}
}

\annot{24}{Questions 6 et 7}{
  \ecrire[vert][8]{(86,43.5)}{EX. (95\%) : $n = 7$ : ME $= 5{,}80$\\ \phantom{EX. (95\%) :} $n = 28$ : ME $= 2{,}43$}
  \entourer[rouge]{(16.5,18.8)}{7}{2.5}
  \ecrire[violet][8]{(61,6.8)}{$E[\overline{X}] = \mu$ : $\overline{X}$ EST SANS BIAIS POUR $\mu$}
}

% ---------------------------------------------------------------------------
% Cas n grand, marge d'erreur et taille d'échantillon
% ---------------------------------------------------------------------------
\annot{26}{Et si la population}{
  % petite densité asymétrique
  \draw[crayon lisse, orange, domain=0:9, samples=40, smooth, shift={(108,78.8)}]
    plot ({4*\x},{5.5*\x*exp(1-\x)});
  \trait[orange]{(108,78.8) -- (146,78.8)}
  \ecrire[orange][7]{(121,86)}{ASYMÉTRIQUE}
  \entourer[rouge]{(107,53.4)}{7}{2.7}
  \ecrire[rouge][8]{(118,37)}{$z$, PAS $t$ :\\ C'EST LE TLC}
  \fleche[rouge][20]{(118,33.5)}{(104,30.3)}
}

\annot{27}{Exemple 3 : mercure}{
  \entourer[violet]{(124.5,73.3)}{7}{2.6}
  \coche{(106,48.8)}
  % axe 0,35 à 0,52 : x = 112 + (v - 0.35) * 38 / 0.17
  \trait[violet]{(112,27) -- (151,27)}
  \foreach \v/\t in {0.37/0{,}37, 0.42/0{,}42, 0.47/0{,}47}{
    \trait[vert]{({112+(\v-0.35)*38/0.17},26) -- ({112+(\v-0.35)*38/0.17},28)}
    \ecrirec[vert][6.5]{({112+(\v-0.35)*38/0.17},23.5)}{\t}}
  \draw[crayon lisse, vert, line width=1.4pt] (116.5,29.5) -- (138.8,29.5);
  \trait[rouge]{(145.5,25) -- (145.5,34)}
  \ecrirec[rouge][6.5]{(145.5,23.3)}{0,5}
  \ecrire[rouge][7]{(132,37.5)}{NORME}
  \fleche[rouge][-20]{(141,34.3)}{(145,33)}
}

\annot{28}{Marge d'erreur et erreur-type}{
  \ecrire[orange][7.5]{(135,32.5)}{RARE !}
  \ecrire[orange][7.5]{(135,27)}{EX. 1}
  \ecrire[orange][7.5]{(135,21.5)}{EX. 3}
  \ecrire[vert][8]{(74,10)}{$\sigma \uparrow \Rightarrow$ ME $\uparrow$ ; \ $n \uparrow \Rightarrow$ ME $\downarrow$ ; \ $1-\alpha \uparrow \Rightarrow$ ME $\uparrow$}
}

\annot{29}{Effet du niveau de confiance}{
  \ecrire[rouge][7.5]{(131,47)}{$n \times 4$ :\\ ME $\div 2$\\ ENVIRON}
  \ecrire[vert][9]{(10,14)}{ME $= z_{\alpha/2}\,\dfrac{s}{\sqrt{n}}$ \ : \ $n \times 4 \Rightarrow \sqrt{n} \times 2 \Rightarrow$ ME $\div\, 2$}
}

\annot{30}{QUIZ}{
  \ecrire[vert][9.5]{(20,62)}{\textcolor{vert}{VRAI} : $\sigma \uparrow \Rightarrow$ ME $= z_{\alpha/2}\,\sigma/\sqrt{n}$ $\uparrow$}
  \ecrire[vert][9.5]{(20,44.5)}{\textcolor{rouge}{FAUX} : $\sqrt{n}$ AU DÉNOMINATEUR $\Rightarrow$ ME $\downarrow$ (PLUS PRÉCIS)}
  \ecrire[vert][9.5]{(20,28)}{\textcolor{vert}{VRAI} : QUANTILE $\uparrow$ : \ $1{,}645 \rightarrow 1{,}96 \rightarrow 2{,}576$}
  \ecrire[vert][8.5]{(45,12.5)}{\textcolor{rouge}{FAUX} : LA TAILLE DE LA POPULATION N'APPARAÎT PAS\\ DANS ME (TANT QU'ELLE EST GRANDE DEVANT $n$)}
}

\annot{31}{Planifier une étude}{
  \ecrire[violet][9]{(116,67.5)}{ON ISOLE $n$}
  \ecrire[vert][9]{(45,13)}{EX. : $n \ge 61{,}01 \Rightarrow n = 62$\\ (AVEC 61, LA MARGE DÉPASSERAIT $E$)}
}

\annot{32}{Exemple 1 (suite) : taille}{
  \ecrire[rouge][8]{(10,55)}{95\% $\Rightarrow$\\ $z_{0{,}025}$}
  \fleche[rouge][-25]{(20,48)}{(48.5,48.5)}
  \ecrire[rouge][7.5]{(121,53.5)}{ARRONDI\\ VERS LE HAUT !}
  \ecrire[vert][8.5]{(121,28.5)}{$= 4 \times 61{,}47$}
}

% ---------------------------------------------------------------------------
% Intervalle de confiance pour une proportion
% ---------------------------------------------------------------------------
\annot{34}{IC pour une proportion}{
  \ecrire[vert][8.5]{(62,56.8)}{$\dfrac{\widehat{P} - p}{\sqrt{p(1-p)/n}} \approx N(0,1)$}
  \ecrire[violet][8]{(119,36.5)}{ERREUR-TYPE\\ ESTIMÉE}
  \fleche[violet][20]{(121,30)}{(114,27)}
  \souligner[rouge]{(43,18.3)}{(104,18.3)}
  \ecrire[rouge][7.5]{(106,23)}{(NB DE SUCCÈS\\ ET D'ÉCHECS)}
}

\annot{35}{Exemple 2 (suite) : le sondage}{
  \ecrire[vert][7.5]{(14,41.9)}{$\Rightarrow$ ENTRE 18,5\% ET 23,5\% DES QUÉBÉCOIS (MÉTHODE FIABLE 19 FOIS SUR 20)}
  \entourer[rouge]{(141.5,26)}{4}{2.4}
  \ecrire[violet][8]{(56,16.5)}{MAX DE $\hat{p}(1-\hat{p})$ : $0{,}5 \times 0{,}5 = 0{,}25$}
  \ecrire[vert][8]{(10,7)}{ICI : $2{,}5\% < 3{,}1\%$ CAR $\hat{p} = 0{,}21$ EST LOIN DE $0{,}5$}
}

\annot{36}{Exemple 4 : prévalence}{
  \ecrire[vert][8.5]{(20,54.3)}{$\hat{p} = 12/80 = 0{,}15$ \ (15\%)}
  \ecrire[vert][8.5]{(20,44.6)}{$n\hat{p} = 12 \ge 5$ \ ET \ $n(1-\hat{p}) = 68 \ge 5$ : OUI}
  \coche{(103,42.5)}
  \ecrire[vert][8.5]{(20,34.4)}{$0{,}15 \pm 1{,}96\sqrt{0{,}15 \times 0{,}85/80} = 0{,}15 \pm 0{,}078 = [0{,}072 \,;\, 0{,}228]$}
  \ecrire[rouge][8.5]{(20,18.8)}{$5\% < 7{,}2\%$ (BORNE INF.) : PEU PLAUSIBLE}
  \ecrire[violet][8]{(66,13)}{DÉTAILS : PAGE SUIVANTE}
}
\apres{36}{
  \ecrire[violet][13]{(7,86)}{EXEMPLE 4 : 12 VEAUX INFECTÉS SUR $n = 80$}
  \ecrire[vert][12]{(7,75)}{1) \ $\hat{p} = \dfrac{12}{80} = 0{,}15$}
  \ecrire[vert][12]{(78,75)}{2) \ $n\hat{p} = 12 \ge 5$\\[1mm] \phantom{2) \ }$n(1-\hat{p}) = 68 \ge 5$}
  \coche{(136,67.5)}
  \ecrire[vert][12]{(7,58)}{3) \ $\widehat{ET} = \sqrt{\dfrac{0{,}15 \times 0{,}85}{80}} = 0{,}0399$}
  \ecrire[vert][12]{(7,44)}{\phantom{3) \ }ME $= 1{,}96 \times 0{,}0399 = 0{,}078$}
  \ecrire[vert][12]{(7,35)}{\phantom{3) \ }IC À 95\% : $0{,}15 \pm 0{,}078 = [0{,}072 \,;\, 0{,}228]$}
  \cadre[rouge]{63.5}{27.3}{107.5}{35.8}
  \ecrire[violet][10]{(7,24)}{LA PRÉVALENCE EST ENTRE 7,2\% ET 22,8\% ;\\ MÉTHODE QUI CAPTE LE VRAI $p$ 95 FOIS SUR 100.}
  % axe des proportions : x = 12 + v * 400
  \begin{scope}
    \trait[violet]{(95,8) -- (152,8)}
    \foreach \v/\t in {0/0, 0.1/0{,}10, 0.2/0{,}20}{
      \trait[violet]{({95+\v*180},7) -- ({95+\v*180},9)}
      \ecrirec[violet][7]{({95+\v*180},4.5)}{\t}}
    \draw[crayon lisse, vert, line width=1.6pt] ({95+0.072*180},11) -- ({95+0.228*180},11);
    \ecrirec[vert][7]{({95+0.15*180},14.5)}{IC}
    \trait[rouge]{({95+0.05*180},6) -- ({95+0.05*180},15)}
    \ecrire[rouge][8]{(96,21)}{4) 5\% : HORS DE L'IC}
    \fleche[rouge][-20]{(100,16.5)}{({95+0.05*180+0.5},14.5)}
  \end{scope}
}

\annot{37}{Taille d'échantillon pour une proportion}{
  \ecrire[violet][8.5]{(98,74)}{$E$ = MARGE D'ERREUR VOULUE}
  % parabole p(1-p)
  \begin{scope}[shift={(116,48.3)}]
    \trait[violet]{(0,0) -- (30,0)}
    \draw[crayon lisse, orange, domain=0:1, samples=40, smooth] plot ({30*\x},{28*\x*(1-\x)});
    \draw[orange, dashed] (15,0) -- (15,7);
    \ecrirec[orange][6.5]{(8.5,7.6)}{0,25}
    \ecrire[orange][6.5]{(18,9.4)}{$p(1-p)$}
    \ecrire[violet][6]{(15.8,3.2)}{0,5}
  \end{scope}
  \ecrire[rouge][8]{(120,15.5)}{$\leftarrow$ PRUDENT}
  \ecrire[vert][8]{(111,7.5)}{$\leftarrow$ 139 TIQUES\\ \ \ \ DE MOINS}
}

\annot{38}{Avec R : prop.test}{
  \ecrire[vert][8]{(72,64.5)}{$\approx$ FORMULE : $[0{,}185 \,;\, 0{,}235]$}
  \ecrire[violet][8]{(72,55)}{VEAUX : BINOM.TEST(12, 80) $\rightarrow [0{,}080 \,;\, 0{,}247]$\\ \ \ (FORMULE : $[0{,}072 \,;\, 0{,}228]$)}
}

\annot{39}{QUIZ}{
  \ecrire[vert][8.5]{(20,58)}{\textcolor{rouge}{FAUX} : ME $= 1{,}96\sqrt{\hat{p}(1-\hat{p})/n}$ DÉPEND DE $n$,\\ PAS DE LA TAILLE DE LA POPULATION}
  \ecrire[vert][8.5]{(20,36.5)}{\textcolor{vert}{VRAI} : $0{,}1 \times 0{,}9 = 0{,}09 < 0{,}25$ ; \ ME $= 1{,}9\%$ VS $3{,}1\%$}
  \ecrire[vert][8.5]{(20,15)}{\textcolor{rouge}{FAUX} : $n\hat{p} = 20 \times 0{,}05 = 1 < 5$ $\Rightarrow$ BINOM.TEST(1, 20)}
}

% ---------------------------------------------------------------------------
% Intervalle de confiance pour sigma^2
% ---------------------------------------------------------------------------
\annot{41}{Pourquoi estimer une variance}{
  \entourer[violet]{(44,73.4)}{8.5}{2.6}
  \ecrire[rouge][8]{(106,17.5)}{$\leftarrow$ KHI-DEUX\\ \ \ \ À $n-1$ D.L.}
}

\annot{42}{Quantiles de la loi du khi-deux}{
  \ecrire[rouge][7.5]{(57,67)}{ASYMÉTRIQUE :\\ PAS DE $\pm$ !}
  \ecrire[violet][9]{(10,24.5)}{TABLE : $P(\chi^2_k > \chi^2_{k;\gamma}) = \gamma$}
  \ecrirec[bleu][8]{(15,15.5)}{$k$}
  \ecrirec[bleu][8]{(40,15.5)}{0,975}
  \ecrirec[bleu][8]{(62,15.5)}{0,025}
  \trait[bleu]{(8,12.5) -- (72,12.5)}
  \trait[bleu]{(22,17.5) -- (22,2)}
  \ecrirec[bleu][8]{(15,9)}{6}
  \ecrirec[bleu][8]{(40,9)}{1,237}
  \ecrirec[bleu][8]{(62,9)}{14,449}
  \ecrirec[bleu][8]{(15,4)}{19}
  \ecrirec[bleu][8]{(40,4)}{8,907}
  \ecrirec[bleu][8]{(62,4)}{32,852}
  \surligner[rose]{(9,6.8)}{(71,11.2)}
  \ecrire[vert][8]{(78,11)}{$n = 7$, 95\% : LIGNE 6\\ (EXEMPLE 6 : LIGNE 19)}
}

\annot{43}{Construction de l'IC pour}{
  \ecrire[rouge][8]{(96,83)}{$\sim \chi^2_{n-1}$ : LOI CONNUE}
  \fleche[rouge][20]{(98,79)}{(90,72.5)}
  \entourer[orange]{(41.5,31.5)}{7}{2.8}
  \ecrire[violet][8.5]{(40,12)}{$\chi^2_{n-1;\alpha/2}$ = LE GRAND QUANTILE, AU DÉNOMINATEUR\\ $\Rightarrow$ PETITE VALEUR $\Rightarrow$ BORNE INFÉRIEURE}
  \fleche[violet][20]{(44,12.5)}{(42,28.6)}
}
\apres{43}{
  \ecrire[violet][13]{(7,86)}{ON ISOLE $\sigma^2$ :}
  \ecrire[vert][11]{(7,76)}{$\chi^2_{n-1;1-\alpha/2} \le \dfrac{(n-1)S^2}{\sigma^2} \le \chi^2_{n-1;\alpha/2}$}
  \ecrire[violet][9]{(7,62)}{ON INVERSE (TOUT EST POSITIF) : \textcolor{rouge}{LES $\le$ CHANGENT DE SENS}}
  \ecrire[vert][11]{(7,55)}{$\dfrac{1}{\chi^2_{n-1;1-\alpha/2}} \ge \dfrac{\sigma^2}{(n-1)S^2} \ge \dfrac{1}{\chi^2_{n-1;\alpha/2}}$}
  \ecrire[violet][9]{(7,40)}{ON MULTIPLIE PAR $(n-1)S^2$ :}
  \ecrire[vert][11]{(7,33)}{$\dfrac{(n-1)S^2}{\chi^2_{n-1;\alpha/2}} \le \sigma^2 \le \dfrac{(n-1)S^2}{\chi^2_{n-1;1-\alpha/2}}$}
  \cadre[rouge]{5}{17}{86}{34}
  \ecrire[orange][9]{(7,13)}{GRAND QUANTILE $\rightarrow$ BORNE INF.\\ PETIT QUANTILE $\rightarrow$ BORNE SUP.}
  % densité khi-deux à 6 d.l. : x = 2 mm par unité, y = 185 mm par unité
  \begin{scope}[shift={(106,42)}]
    \fill[pattern=north east lines, pattern color=vert] (0,0) --
      plot[domain=0.01:1.237, samples=20] ({2*\x},{11.5625*\x*\x*exp(-\x/2)}) -- (2.47,0) -- cycle;
    \fill[pattern=north east lines, pattern color=vert] (28.9,0) --
      plot[domain=14.449:23, samples=20] ({2*\x},{11.5625*\x*\x*exp(-\x/2)}) -- (46,0) -- cycle;
    \draw[crayon lisse, violet, line width=1pt] plot[domain=0:23, samples=90, smooth] ({2*\x},{11.5625*\x*\x*exp(-\x/2)});
    \trait[violet]{(-1,0) -- (48,0)}
    \trait[violet]{(0,-1) -- (0,30)}
    \ecrire[violet][10]{(20,31)}{$\chi^2_{n-1}$}
    \ecrirec[vert][8]{(1.5,-4.5)}{$\chi^2_{n-1;1-\frac{\alpha}{2}}$}
    \ecrirec[vert][7]{(1.5,-10)}{(PETIT)}
    \ecrirec[vert][8]{(29,-4.5)}{$\chi^2_{n-1;\frac{\alpha}{2}}$}
    \ecrirec[vert][7]{(29,-10)}{(GRAND)}
    \ecrirec[vert][9]{(-6,14)}{$\frac{\alpha}{2}$}
    \fleche[vert][-20]{(-5,10.5)}{(1.4,1.5)}
    \ecrirec[vert][9]{(41,13)}{$\frac{\alpha}{2}$}
    \fleche[vert][20]{(40,9.5)}{(33,2)}
    \ecrirec[violet][9]{(15,8)}{$1-\alpha$}
  \end{scope}
}

\annot{44}{Exemple 1 (suite) : variabilité}{
  \ecrire[violet][7.5]{(131,74.5)}{$= 6 \times 39{,}33$}
  % axe 0 à 200 : x = 119 + v * 0.155
  \trait[violet]{(118,55) -- (151,55)}
  \draw[crayon lisse, vert, line width=1.5pt] ({119+16.33*0.155},57) -- ({119+190.73*0.155},57);
  \trait[rouge]{({119+39.33*0.155},53.5) -- ({119+39.33*0.155},59.5)}
  \ecrirec[rouge][7]{({119+39.33*0.155},61.5)}{$s^2$}
  \ecrire[rouge][7]{(128,63.5)}{PAS AU CENTRE !}
}

\annot{45}{Exemple 6 : homogénéité}{
  \ecrire[vert][8.5]{(20,48.2)}{$(n-1)s^2 = 19 \times 1{,}6^2 = 48{,}64$ \ ; \ $\sigma^2 \in [\,1{,}48 \,;\, 5{,}46]$ G$^2$}
  \ecrire[vert][8.5]{(20,37.7)}{$\sigma \in [\sqrt{1{,}48} \,;\, \sqrt{5{,}46}] = [\,1{,}22 \,;\, 2{,}34]$ G}
  \ecrire[rouge][8.5]{(20,27.5)}{$1 \notin [\,1{,}22 \,;\, 2{,}34]$ $\Rightarrow$ PEU PLAUSIBLE}
  \ecrire[violet][8]{(100,12)}{DÉTAILS : PAGE SUIVANTE}
}
\apres{45}{
  \ecrire[violet][13]{(7,86)}{EXEMPLE 6 : $n = 20$, \ $s = 1{,}6$ G, \ IC À 95\%}
  \ecrire[vert][12]{(7,76)}{$(n-1)s^2 = 19 \times 1{,}6^2 = 19 \times 2{,}56 = 48{,}64$}
  \ecrire[vert][12]{(7,65)}{1) \ $\sigma^2 \in \left[\dfrac{48{,}64}{32{,}852} \,;\, \dfrac{48{,}64}{8{,}907}\right] = [\,1{,}48 \,;\, 5{,}46]$ G$^2$}
  \ecrire[orange][8]{(30,49)}{$\uparrow$ GRAND QUANTILE\\ \ \ \ $\rightarrow$ BORNE INF.}
  \ecrire[vert][12]{(7,39)}{2) \ $\sigma \in \left[\sqrt{1{,}48} \,;\, \sqrt{5{,}46}\right] = [\,1{,}22 \,;\, 2{,}34]$ G}
  \ecrire[vert][12]{(7,26)}{3) \ $\sigma = 1$ G : HORS DE L'IC}
  \ecrire[rouge][11]{(7,17)}{$\Rightarrow$ PEU PLAUSIBLE : LES SOURIS\\ SONT PLUS VARIABLES QUE PRÉTENDU}
  \ecrire[violet][9]{(100,36)}{NON SYMÉTRIQUE AUTOUR\\ DE $s = 1{,}6$ (CENTRE : 1,78)}
  % axe 0 à 3 : x = 96 + v * 18
  \trait[violet]{(95,12) -- (153,12)}
  \foreach \v in {0,1,2,3}{
    \trait[violet]{({96+\v*18},11) -- ({96+\v*18},13)}
    \ecrirec[violet][8]{({96+\v*18},8)}{\v}}
  \draw[crayon lisse, vert, line width=1.6pt] ({96+1.22*18},15) -- ({96+2.34*18},15);
  \ecrirec[vert][8]{({96+1.78*18},19)}{IC POUR $\sigma$}
  \trait[rouge]{(114,10) -- (114,19)}
  \ecrirec[rouge][8]{(111,22)}{1 G}
}

\annot{46}{QUIZ}{
  \ecrire[vert][9]{(20,66)}{\textcolor{rouge}{FAUX} : NON SYMÉTRIQUE AUTOUR DE $s^2$ \ (16,33 ; 39,33 ; 190,73)}
  \ecrire[vert][9]{(20,46.5)}{\textcolor{vert}{VRAI} : $\sqrt{\ }$ CROISSANTE : $P(a \le \sigma^2 \le b) = P(\sqrt{a} \le \sigma \le \sqrt{b})$}
  \ecrire[vert][9]{(20,25)}{\textcolor{rouge}{FAUX} : PAS DE TLC POUR $S^2$ : LA MÉTHODE EXIGE LA\\ NORMALITÉ, MÊME SI $n$ EST GRAND}
}

% ---------------------------------------------------------------------------
% Synthèse
% ---------------------------------------------------------------------------
\annot{48}{Tableau récapitulatif}{
  \ecrire[rouge][8]{(118,59.5)}{$\leftarrow$ LE PLUS FRÉQUENT}
  \ecrire[rouge][8]{(118,37.5)}{$\leftarrow$ PAS DE $\pm$ !}
  \surligner[jaune]{(12,53)}{(113,57.5)}
}

\produire{Chapitre_3b}
\end{document}
